\section{Study Design}\label{sec:design}

We ask four research questions.
\textbf{RQ1} asks which dependence definitions give sound and dependence-complete slices.
\textbf{RQ2} asks how much smaller than the wire slice a slice can safely be, and under which inputs the
extra precision holds.
\textbf{RQ3} asks what each definition does to fault localization, by hand, with mutation-based fault
localization and with language models.
\textbf{RQ4} asks what each definition costs in time and memory.

\paragraph*{Pre-registration}
The protocol, including the slicers, metrics, tests, thresholds and the decision rules for the
paper's framing, was committed before any run. Every later change is recorded as a dated deviation
with its reason, and those made after a result are marked as such. The changes that affect this paper
add baselines. The full classical PDG and an order-preserving mapping for Qiskit's \texttt{LightCone}
pass were added after the first runs, and both are reported beside the pre-registered results. No
result enforces an expected outcome. The replication package contains the protocol, the deviations,
the logs of every run, and the commit each run used.

\paragraph*{Circuits}
Corpus~A has \numNCircuits{} circuits of at most \numMaxQubits{} qubits from three benchmark suites,
MQT~Bench (\numCircuitsMqt{} circuits)~\cite{quetschlich2023mqt}, QASMBench
(\numCircuitsQasm{})~\cite{li2023qasmbench} and VeriQBench (\numCircuitsVeriq{})~\cite{chen2022veriqbench}.
Every qubit of every circuit is a criterion, which gives \numNPairs{} circuit--qubit pairs, of which
\numNDynamicPairs{} are in circuits with mid-circuit measurement, reset or classical feedback. We also
generate \numFamCircuits{} circuits with known structure. They cover depth sweeps at fixed width,
multipartite states that are partly uncomputed, chains that entangle, disentangle and re-entangle,
indirect entanglement through a third qubit, and inverse pairs $U \dots U^{\dagger}$. For scale,
corpus~L has \numCorpusLCircuits{} larger circuits of 15 to 20 qubits.

\paragraph*{Oracle and ground truth}
All slices are judged by one independent oracle. It simulates the circuit's unconditional state exactly,
by branching over measurement outcomes, and shares no code with any slicer. It agrees with Qiskit's
statevector simulation to within \numOracleErr{} on random circuits. A slice preserves the criterion when
the Uhlmann fidelity of the criterion's reduced state is at least $1 - 10^{-6}$. The dependences of
Definition~\ref{def:dep} are measured by inserting \numInfluencePerturbations{} random unitaries after
every operation and comparing the criterion's state at \numInfluenceInputs{} inputs, $\ket{0\dots0}$ and
ten random product states. A dependence catalogue of \numCatEntries{} hand-derived cases, including the
examples of Section~\ref{sec:dependence}, checks the ground truth and every slicer before the runs.

\paragraph*{Slicers}
Each rule of Section~\ref{sec:dependence} has an implementation. The classical PDG and def-use
implement their read and write rules directly. The wire rule is a backward slice over per-qubit data edges
and measurement-to-condition edges. Entanglement closure adds the coupling-component edges of our
documented reading of Yu and Zhao's graph~\cite{yu2025qpdg}, with every interpretation choice logged
before the runs. Commutation uses Qiskit's \texttt{LightCone} pass, and state-based pruning uses the
algorithm and thresholds of Section~\ref{sec:dependence}. A variant of state-based pruning takes a set of
qubits as its criterion.

\paragraph*{Robustness}
A slice computed at $\ket{0\dots0}$ is checked again, with its positions fixed, under
\numRobDraws{} random computational-basis inputs, \numRobDraws{} random product inputs and, for circuits
with rotation angles, \numRobDraws{} random rebindings of every angle. A pair is robust if the slice
preserves the criterion in every draw.

\paragraph*{Faults}
We seed faults with \numMutOperators{} operators taken from Muskit~\cite{mendiluze2021muskit},
QMutPy~\cite{fortunato2022qmutpy} and the fault model of García de la Barrera et
al.~\cite{garcia2022testing}. They insert, delete and replace gates, insert and delete mid-circuit
measurements, swap a control with a target, remove a control, re-target an operand, corrupt the
preparation of an entangling gate, and perturb an angle. Of the \numMutants{} mutants,
\numMutantsKilled{} change the measured distribution of at least one qubit. Each such qubit is a unit,
which gives \numRqThreeUnits{} units whose criterion is a qubit the developer would see is wrong. A further
\numMutantsInjured{} mutants change only a qubit's phase and form \numRqThreeSecondary{} secondary units.
For each unit we record whether the slice of the faulty circuit contains the fault site and how many
instructions it keeps. The real bugs come from Bugs4Q~\cite{zhao2023bugs4q}. Of \numBugsFourQTriaged{}
entries, \numBugsFourQIncluded{} are circuit-level faults whose buggy and fixed programs can be ported
mechanically to the current Qiskit. Of these, \numBugsFourQKilled{} change a qubit's measured
distribution, \numBugsFourQInjured{} changes only a phase, \numBugsFourQAlive{} changes no qubit, and
\numBugsFourQExcluded{} exceed the oracle. We map their slices back to the lines of the original Qiskit
programs.

\paragraph*{Fault localization}
Mutation-based fault localization ranks instructions with the Ochiai formula, as for quantum programs
in Ishimoto et al.~\cite{ishimoto2025mbfl}, over the whole circuit and within each slice. In the language
model task, a model sees a faulty circuit, or one of its slices, and the wrong qubit, and names up to
three suspect instructions. We score whether the fault site is among them, on \numTTwoItems{} units
sampled across operators, with two Claude models and two open-weight Qwen2.5-Coder models.

\paragraph*{Statistics}
The confirmatory hypotheses H1 to H5 are tested with exact McNemar tests on paired binary outcomes and
Wilcoxon signed-rank tests on paired sizes, with Holm correction across the five. Effect sizes come
first, with bootstrap 95\% confidence intervals (10{,}000 resamples) and Wilson intervals for
proportions. Everything else is descriptive.
